% Generated by scripts/figures/make_paper_figures.py; do not edit. Needs \usepackage{booktabs,graphicx,xcolor}.
\begin{table*}[t]
\centering
\caption{X-M: one LLM for detection and instance segmentation (final report set, Qwen3.5-9B, F1 $\times$100). Category-level detection and instance segmentation with Qwen3.5-9B on the final report set: 1,000 COCO val2017 images, 2,940 (image, category) pairs over 79 categories, disjoint by image from the tuning and decision sets; X-M's training images and all evaluation images are disjoint. Only the categories present in an image are asked (positive only, as in Rex-Omni's evaluation). F1 $\times$100 = Rex-Omni's class-macro F1 (preregistration XM-3a): COCOeval-style greedy matching per (image, category) in output order, precision and recall per category over the set, averaged over categories, F1 = 2PR/(P+R); IoU 0.5 on boxes (detection) or masks (instance segmentation). P / R: detection precision and recall at IoU 0.5. F1@mIoU: P and R each averaged over IoU 0.5:0.05:0.95 before the harmonic mean. Per-row CIs: 95\% bootstrap over images from each run (1,000 resamples). Main $-$ row: X-M's main setting minus the row (for detection, X-M's boxes minus the row's boxes), paired bootstrap over images with the macro average recomputed in every resample, 95\%, 10,000 resamples (seed 20260929). -- = not applicable (zero-shot boxes and Rex-Omni give boxes only; no row is paired with itself). The X-M rows share call 1's boxes and the two base rows share the base model's boxes, so their detection columns are identical. No entry is in bold. On the final report set, the spline rasterised directly, without snapping to image edges, is only 1.2 points [+0.6, +1.7] below the main setting: the most direct evidence that no dedicated mask decoder is needed. The shape comes entirely from the LLM, and GrabCut adds only about 1 point. External reference: Rex-Omni (IDEA-Research/Rex-Omni at 0e5693d), a detector built on Qwen2.5-VL-3B, i.e. a learned vision model; run by us separately with its official prompts, preprocessing and parsing (deviations recorded with the run, D178), detection only; its outputs are used only for scoring, never in our data or training. Joint = its own protocol: one query per image listing all the image's positive categories. X-M $-$ Rex-Omni (joint): $-$1.6 [$-$3.3, +0.5] points; the interval includes 0, so X-M and Rex-Omni are comparable on these images. Gray rows: detection F1@0.5 (Rex-Omni's definition: class-macro precision and recall, F1 = 2PR/(P+R)) on the full COCO val2017 (4,952 images) as reported in the literature, not on our 1,000-image subset; reference only, never paired; sources under the table. Data: RUN-0109 (X-M, base; hpcgpu07); RUN-0111 (Rex-Omni; hpcgpu06); X-M and the base model on the same node (D93-1); Rex-Omni, an external reference, on hpcgpu06 (D180); readings and paired differences from xm\_results.json (computed in RUN-0113 on CPU); literature from xm\_sota.json. Definitions of the settings are under the table.}
\label{tab:xm-main}
\footnotesize
\setlength{\tabcolsep}{3pt}
\resizebox{\textwidth}{!}{%
\begin{tabular}{lcccccc}
\toprule
System & Det. F1@0.5 [95\% CI] & Det. P / R & Det. F1@mIoU & Inst. seg. F1@0.5 [95\% CI] & Main $-$ row: det. F1 [95\% CI] & Main $-$ row: inst. seg. F1 [95\% CI] \\
\midrule
\multicolumn{7}{l}{\emph{Qwen3.5-9B base model, zero-shot (same prompts and post-processing)}} \\
Base, zero-shot boxes & $64.2$\,{\scriptsize $[62.9, 67.7]$} & $62.6$ / $66.0$ & $47.7$ & -- & $+5.3$\,{\scriptsize $[+3.5, +6.5]$} & -- \\
Base boxes → GrabCut & $64.2$\,{\scriptsize $[62.9, 67.7]$} & $62.6$ / $66.0$ & $47.7$ & $39.3$\,{\scriptsize $[37.6, 42.1]$} & -- & $+19.2$\,{\scriptsize $[+17.1, +21.4]$} \\
\midrule
\multicolumn{7}{l}{\emph{X-M (ours): one LoRA on Qwen3.5-9B}} \\
X-M: spline → GrabCut (main setting) & $69.5$\,{\scriptsize $[68.2, 72.5]$} & $71.8$ / $67.4$ & $52.1$ & $58.5$\,{\scriptsize $[56.9, 61.5]$} & -- & -- \\
\midrule
\multicolumn{7}{l}{\emph{Ablations of X-M: the same boxes, other masks (paired against the main setting)}} \\
X-M: spline rasterised directly, no snapping & $69.5$\,{\scriptsize $[68.2, 72.5]$} & $71.8$ / $67.4$ & $52.1$ & $57.3$\,{\scriptsize $[55.7, 60.5]$} & -- & $+1.2$\,{\scriptsize $[+0.6, +1.7]$} \\
X-M: in-box ellipse → GrabCut & $69.5$\,{\scriptsize $[68.2, 72.5]$} & $71.8$ / $67.4$ & $52.1$ & $51.2$\,{\scriptsize $[50.1, 54.2]$} & -- & $+7.3$\,{\scriptsize $[+5.1, +8.8]$} \\
X-M: boxes → GrabCut (no spline) & $69.5$\,{\scriptsize $[68.2, 72.5]$} & $71.8$ / $67.4$ & $52.1$ & $41.2$\,{\scriptsize $[39.5, 43.9]$} & -- & $+17.3$\,{\scriptsize $[+15.4, +19.6]$} \\
\midrule
\multicolumn{7}{l}{\emph{External reference: Rex-Omni, a learned detector run by us (detection only)}} \\
Rex-Omni, joint (its own protocol) & $71.1$\,{\scriptsize $[69.8, 73.4]$} & $76.3$ / $66.6$ & $52.0$ & -- & $-1.6$\,{\scriptsize $[-3.3, +0.5]$} & -- \\
\midrule
\multicolumn{7}{l}{\textcolor{gray}{\emph{Literature: COCO val2017, all 4,952 images, as reported (not the same images; reference only, not paired)}}} \\
\textcolor{gray}{Qwen2.5-VL-3B, MLLM} & \textcolor{gray}{$64.7$} & \textcolor{gray}{--} & \textcolor{gray}{--} & \textcolor{gray}{--} & \textcolor{gray}{--} & \textcolor{gray}{--} \\
\textcolor{gray}{Qwen2.5-VL-7B, MLLM} & \textcolor{gray}{$64.0$} & \textcolor{gray}{--} & \textcolor{gray}{--} & \textcolor{gray}{--} & \textcolor{gray}{--} & \textcolor{gray}{--} \\
\textcolor{gray}{SEED1.5-VL, MLLM} & \textcolor{gray}{$71.3$} & \textcolor{gray}{--} & \textcolor{gray}{--} & \textcolor{gray}{--} & \textcolor{gray}{--} & \textcolor{gray}{--} \\
\textcolor{gray}{Rex-Omni, MLLM detector} & \textcolor{gray}{$72.0$} & \textcolor{gray}{--} & \textcolor{gray}{--} & \textcolor{gray}{--} & \textcolor{gray}{--} & \textcolor{gray}{--} \\
\textcolor{gray}{Grounding DINO-T, specialist detector} & \textcolor{gray}{$69.8$} & \textcolor{gray}{--} & \textcolor{gray}{--} & \textcolor{gray}{--} & \textcolor{gray}{--} & \textcolor{gray}{--} \\
\textcolor{gray}{DINO-Swin-L, specialist detector} & \textcolor{gray}{$75.6$} & \textcolor{gray}{--} & \textcolor{gray}{--} & \textcolor{gray}{--} & \textcolor{gray}{--} & \textcolor{gray}{--} \\
\bottomrule
\end{tabular}}
\par\smallskip
\begin{minipage}{\textwidth}\footnotesize
X-M (ours): one LoRA on Qwen3.5-9B, trained with X-2's recipe plus two COCO train2017 streams (listing a category's boxes; one spline per given box); supervision only from ground truth and self-distillation replay of the base model; no learned vision model in training or inference. Per (image, category) two calls, single pass, no retries: call 1 lists the category's boxes (at most 30); call 2 traces each box's outline as a closed B-spline with 16 control points. Deterministic post-processing: NMS at IoU 0.7 over the listed boxes, control points clipped to their box. Main setting: the rasterised spline seeds GrabCut as a trimap for each instance (snapping to image edges); GrabCut (OpenCV graph cut) has no learned parameters. Ablations with the same boxes: spline rasterised directly (the spline's interior is the mask, no GrabCut); in-box ellipse → GrabCut (the spline replaced by the ellipse inscribed in the box); boxes → GrabCut (the box alone seeds GrabCut, no spline). Base: the zero-shot Qwen3.5-9B with the same prompts (xm-prompts-v1, frozen) and post-processing. Final report set: 436 of 7,591 predicted instances (5.7\%) fall back to boxes → GrabCut in the main setting, because call 2 returned a different number of splines than boxes (D168). Literature rows: Qwen2.5-VL-3B (MLLM, one category per query as in our protocol): Rex-Omni, arXiv 2510.12798 ('Detect Anything via Next Point Prediction'), Table 2, COCO val2017 detection F1@0.5. Qwen2.5-VL-7B (MLLM, one category per query as in our protocol): Rex-Omni, arXiv 2510.12798, Table 2, COCO val2017 detection F1@0.5. SEED1.5-VL (MLLM, not one category per query): Rex-Omni, arXiv 2510.12798, Table 2, COCO val2017 detection F1@0.5. Rex-Omni (MLLM detector built on Qwen2.5-VL-3B, its own protocol): Rex-Omni, arXiv 2510.12798, Table 2, COCO val2017 detection F1@0.5; we also run its released weights (IDEA-Research/Rex-Omni) on our final report set (the non-gray Rex-Omni rows). Grounding DINO-T (specialist detector at the best score threshold of a 0-1 sweep): Rex-Omni, arXiv 2510.12798, Table 2, COCO val2017 detection F1@0.5. DINO-Swin-L (specialist detector at the best score threshold of a 0-1 sweep): Rex-Omni, arXiv 2510.12798, Table 2, COCO val2017 detection F1@0.5.
\end{minipage}
\end{table*}
